\documentclass[12pt,a4paper]{article}

% --- Настройка языка и шрифтов для LuaLaTeX ---
\usepackage{fontspec}
\defaultfontfeatures{Ligatures=TeX,Scale=MatchLowercase}

\setmainfont{Alice-Regular}[
    Path=font/,
    Extension=.ttf,
    BoldFont = Alice-Regular,
    ItalicFont = Alice-Regular,
    BoldItalicFont = Alice-Regular,
    BoldFeatures = {FakeBold=1.6},
    ItalicFeatures = {FakeSlant=0.2},
    BoldItalicFeatures = {FakeBold=1.6,FakeSlant=0.2}
]

\newfontfamily\cyrillicfont{Alice-Regular}[
    Path=font/,
    Extension=.ttf,
    BoldFont = Alice-Regular,
    ItalicFont = Alice-Regular,
    BoldItalicFont = Alice-Regular,
    BoldFeatures = {FakeBold=1.6},
    ItalicFeatures = {FakeSlant=0.2},
    BoldItalicFeatures = {FakeBold=1.6,FakeSlant=0.2}
]

\usepackage{polyglossia}
\setmainlanguage{russian}
\setsansfont{TeX Gyre Heros}
\setmonofont{TeX Gyre Cursor}

\usepackage{graphicx}        					% Графика
\usepackage{amsmath, amsfonts, amssymb, amsthm}	% Математика
\usepackage{mathtools}
\usepackage{chemfig}
\usepackage[version=4,arrows=pgf-filled,
            textfontname=sffamily,
            mathfontname=mathsf]{mhchem}
\usepackage[protrusion=false]{microtype}       % Микротипографика
\usepackage{braket} 						% Braket нотация для квантовой механики
\usepackage{unicode-math} 					% Современный пакет для математики в LuaLaTeX
\usepackage{longtable}       				% Таблицы с разрывом страниц
\usepackage{tikz}           				% Рисование
\usepackage[europeanresistors]{circuitikz}	% Рисование электрических схем в европейском стиле
\usepackage{listings}						% Для отображения кода
\usepackage{tabularx}       				% Для таблиц во всю ширину
\usepackage{array}		  					% Для расширенных возможностей таблиц
\usepackage{cancel}                % Перечёркивать
\usepackage{float}          				% Для точного позиционирования фигур
\usepackage{pgfplots}                          % Для построения графиков
\usepackage{caption}                           % Для настройки подписей к рисункам и таблицам
\usepackage{xfp}                               % Для вычислений в LaTeX
\usepackage{luacode}                           % Для выполнения Lua-кода в документе    
\pgfplotsset{compat=1.18}
\usepackage[paper=letterpaper,
            top=20mm,
            bottom=20mm,
            left=10mm,
            right=10mm,
            footskip=0.4in 
]{geometry}
\usepackage{colortbl}

\usetikzlibrary{arrows.meta}

% --- Колонтитулы ---
\usepackage{fancyhdr}
\pagestyle{fancy}
\fancyhf{}

% --- Теоремы и окружения ---
\theoremstyle{plain}
\newtheorem{theorem}{Theorem}[section]
\newtheorem{lemma}[theorem]{Lemma}
\newtheorem{proposition}[theorem]{Proposition}
\newtheorem{corollary}[theorem]{Corollary}

\theoremstyle{definition}
\newtheorem{definition}[theorem]{Definition}

\theoremstyle{remark}
\newtheorem{remark}[theorem]{Remark}

\numberwithin{equation}{section}


\renewcommand{\headrulewidth}{0.4pt}
\renewcommand{\footrulewidth}{0.4pt}
\renewcommand\lstlistingname{Исходный код}

\lstdefinestyle{main}{
	backgroundcolor=\color{gray!8},   
	commentstyle=\color{green!50!black},
	keywordstyle=\color{blue!80!black},
	numberstyle=\small\color{darkgray},
	stringstyle=\rmfamily\color{orange!90!black},
	basicstyle=\ttfamily\small,
	breakatwhitespace=false,         
	breaklines=true,                   
	keepspaces=true,                 
	numbers=left,                    
	numbersep=8pt,                  
	showspaces=false,                
	showstringspaces=false,
	showtabs=false,                  
	tabsize=4
}	

\lstset{style=main}

\newcolumntype{C}{>{\centering\arraybackslash}X}
\renewcommand{\arraystretch}{1.25}

\begin{luacode*}
include = require("include")
\end{luacode*}

\begin{document}

\begin{center}
    \Large Лабораторная работа \number 4
\end{center}

\begin{center}
    \large Фотолитография
\end{center}

\textbf{Цель работы}: ознакомиться с процессом фотолитографии; 
исследовать режимы операций контактной фотолитографии.

\section{Теоретические сведения}

Литография — процесс переноса геометрического рисунка шаблона на поверхность 
кремниевой пластины с помощью чувствительных к излучению покрытий.

По типу используемого излучения литографию подразделяют на оптическую 
(фотолитографию, $\lambda = 10 \dots 450$~нм), рентгеновскую ($\lambda = 0{,}4 \dots 1{,}5$~нм), 
ионно-лучевую ($\lambda = 0{,}05 \dots 1$~нм) и электронно-лучевую ($\lambda \approx 0{,}01$~нм).

\subsection{Оптическая литография и дифракционный предел}

Ключевым преимуществом фотолитографии является возможность параллельного и одновременного 
переноса топологического рисунка, содержащего миллионы элементов, на всю площадь кристалла 
или пластины. Это обеспечивает высокую производительность и экономическую эффективность 
производства интегральных схем со степенью интеграции $10^7 \dots 10^{10}$ транзисторов на чип.

По способу передачи изображения различают три варианта реализации оптической печати:
\begin{enumerate}
  \item \textbf{Контактная печать}: фотошаблон механически прижимается непосредственно к слою 
  фоторезиста. Метод обеспечивает высокое разрешение, однако минимальный размер ограничен 
  величиной $D \propto \sqrt{\lambda}$, а непосредственный контакт приводит к быстрому механическому 
  износу и дефектам шаблона.
  \item \textbf{Контактная печать с зазором (проксимити-печать)}: между шаблоном и резистом 
  устанавливается зазор $10 \dots 20$~мкм. Исключается механическое повреждение шаблона, но 
  дифракционные явления на зазоре существенно ухудшают резкость краев и разрешающую способность.
  \item \textbf{Проекционная печать с уменьшением масштаба}: изображение шаблона переносится 
  на пластину высокоапертурной оптической системой с уменьшением (обычно $4:1$ или $5:1$). 
  Шаблон не повреждается, а дефекты и погрешности на нем масштабируются в сторону уменьшения.
\end{enumerate}

Минимальный разрешаемый размер элемента $D$ в проекционной литографии ограничен дифракцией 
и описывается критерием Рэлея — Аббе:
\begin{equation*}
  D = k_1 \frac{\lambda}{\mathrm{NA}},
\end{equation*}
где $\lambda$ — длина волны актиничного излучения; $\mathrm{NA} = n \sin\theta$ — числовая 
апертура объектива ($n$ — показатель преломления среды, $\theta$ — апертурный угол); 
$k_1$ — технологический коэффициент, зависящий от когерентности освещения, свойств фоторезиста 
и способов обработки ($k_1 \sim 0{,}25 \dots 0{,}61$).

Глубина резкости оптической системы ($\mathrm{DOF}$), определяющая допустимую неровность 
подложки, также зависит от длины волны и апертуры:
\begin{equation*}
  \mathrm{DOF} = k_2 \frac{\lambda}{\mathrm{NA}^2}.
\end{equation*}

\subsection{Источники излучения и преодоление дифракционного предела}

Для снижения минимального размера $D$ в оптической литографии используют переход к 
более коротким длинам волн актиничного излучения:
\begin{itemize}
  \item \textbf{Ртутные дуговые лампы}: спектральные линии $g$ ($436$~нм, видимый свет), 
  $h$ ($405$~нм), $i$ ($365$~нм, ближний УФ), DUV ($240 \dots 255$~нм);
  \item \textbf{Эксимерные лазеры}: $\mathrm{KrF}$ ($248$~нм, глубокий УФ), $\mathrm{ArF}$ ($193$~нм, 
  глубокий УФ), $\mathrm{F}_2$ ($157$~нм, вакуумный УФ).
\end{itemize}

Для преодоления дифракционного предела при длине волны $\lambda = 193$~нм в субмикронной 
и нанометровой областях (нормы $45$, $32$ и $22$~нм) разработаны следующие методы:
\begin{enumerate}
  \item \textbf{Коррекция оптической близости (OPC)}: расчетное изменение формы и размеров 
  элементов на фотошаблоне (добавление вспомогательных штрихов, утолщений углов) для 
  компенсации дифракционных искажений волнового фронта.
  \item \textbf{Фазосдвигающие маски (PSM)}: введение фазового сдвига на $180^\circ$ между 
  соседними прозрачными участками шаблона, что за счет интерференционного гашения света 
  на границе формирует нулевую интенсивность и обеспечивает крутой перепад освещенности.
  \item \textbf{Иммерсионная литография}: пространство между последней линзой объектива и 
  пластиной заполняется жидкостью с показателем преломления $n > 1$ (обычно деионизованная 
  очищенная вода с $n \approx 1{,}44$ для $\lambda = 193$~нм). Это повышает эффективную числовую 
  апертуру $\mathrm{NA} > 1$ и эквивалентно уменьшает длину волны в среде $\lambda_\text{эфф} = \lambda / n$.
  \item \textbf{Многократное паттернирование}: формирование плотного рисунка путем последовательного 
  двойного или четверного экспонирования и травления со смещением.
  \item \textbf{Внеосевое освещение (OAI)}: освещение шаблона пучками света под наклоном, при котором 
  в апертуру объектива попадают лучи высших порядков дифракции.
\end{enumerate}

\subsection{Фоторезисты: состав и классификация}

Фоторезисты — сложные светочувствительные полимерные композиции, предназначенные для создания 
защитного рельефа требуемой топологии, устойчивого к травлению в кислотах, щелочах и плазме.

Практически любой фоторезист содержит три основных компонента:
\begin{enumerate}
  \item \textbf{Светочувствительный компонент (фотоактивное соединение)}: вещество, претерпевающее 
  химические превращения под воздействием квантов актиничного света;
  \item \textbf{Пленкообразующая полимерная основа}: полимерная матрица, обеспечивающая адгезию, 
  однородность слоя и механическую/химическую стойкость к травителям;
  \item \textbf{Растворитель}: органическая среда, определяющая вязкость композиции и толщину 
  пленки при центрифугировании.
\end{enumerate}

По характеру изменения растворимости в процессе экспонирования фоторезисты делятся на два класса:
\begin{enumerate}
  \item \textbf{Негативные фоторезисты}: под воздействием излучения растворимость экспонированных 
  участков \textbf{уменьшается}.
  \begin{itemize}
    \item Основой служат эфиры коричной кислоты и поливинилового спирта (поливинилциннамат, ПВЦ) 
    или циклокаучуки с диазидами.
    \item Механизм: при поглощении фотонов рвется двойная связь $\mathrm{C}=\mathrm{C}$ в циннамоильной 
    группе, возникают свободные радикалы, сшивающие линейные полимерные молекулы в пространственную 
    нерастворимую сетку.
    \item Проявляются в органических растворителях (толуол, трихлорэтилен).
    \item Недостаток: в процессе проявления сшитая сетка набухает, что приводит к размытию краев 
    и ограничивает минимальный размер элемента ($l \ge 3 d_\text{пленки}$).
  \end{itemize}
  \item \textbf{Позитивные фоторезисты}: под воздействием излучения растворимость экспонированных 
  участков \textbf{возрастает}.
  \begin{itemize}
    \item Светочувствительным компонентом является нафтохинондиазид (НХД), а пленкообразователем — 
    новолачная фенолоформальдегидная смола.
    \item Проявляются в водных щелочных растворах ($0{,}3 \dots 0{,}5\%$ $\mathrm{KOH}$, тринатрийфосфат).
    \item Преимущество: проявитель не пропитывает неэкспонированный слой, набухание отсутствует, 
    что обеспечивает значительно более высокую разрешающую способность, чем у негативных резистов.
  \end{itemize}
\end{enumerate}

\subsubsection{Фотолиз позитивного фоторезиста на основе НХД}

В исходном состоянии неэкспонированные гидрофобные молекулы НХД распределены по всему объему 
новолачной смолы и служат ингибитором растворения, блокируя доступ щелочного проявителя к смоле. 
Под воздействием УФ-излучения протекает многостадийная реакция:
\begin{enumerate}
  \item При поглощении фотона рвется связь $\mathrm{C}=\mathrm{N}$ и отщепляется молекула азота:
  $\mathrm{N}_2\uparrow$;
  \item Шестичленное углеродное кольцо перестраивается в пятичленное с образованием короткоживущего 
  инденкарбена;
  \item Инденкарбен реагирует с молекулами воды $\mathrm{H}_2\mathrm{O}$, содержащимися в слое резиста, 
  образуя инденкарбоновую кислоту;
  \item Инденкарбоновая кислота легко реагирует с гидроксид-ионами щелочного проявителя с образованием 
  растворимых в воде солей, благодаря чему экспонированные участки быстро вымываются.
\end{enumerate}

\subsection{Основные характеристики фоторезистов}

\subsubsection{Светочувствительность}

Светочувствительность $S$ — величина, обратная экспозиционной дозе $H$, необходимой для 
полного фотохимического превращения резиста:
\begin{equation*}
  S = \frac{1}{H} = \frac{1}{E \cdot t},
\end{equation*}
где $H = E \cdot t$ — экспозиция (Дж/м$^2$ или Вт$\cdot$с/м$^2$), $E$ — облученность поверхности (Вт/м$^2$), 
$t$ — время экспонирования (с). Для позитивных резистов светочувствительность характеризует 
полноту фотолиза НХД, а для негативных — степень задубливания на заданную глубину.

\subsubsection{Разрешающая способность}

Разрешающая способность $R$ — максимальное число параллельных различимых линий равной 
ширины $l$ и зазора между ними $l$, воспроизводимых на $1$~мм длины:
\begin{equation*}
  R = \frac{1}{2l}.
\end{equation*}
Так, для воспроизведения элементов размером $l = 0{,}09$~мкм разрешающая способность резиста 
должна составлять не менее $5555$~лин/мм.

Разрешающая способность позитивных фоторезистов выше негативных по двум причинам:
\begin{itemize}
  \item Отсутствие эффекта набухания полимера в щелочном проявителе;
  \item Отражение света от подложки: в негативном резисте отраженный свет сшивает придонный слой 
  вокруг защитного участка, создавая неустранимый ореол; в позитивном резисте частично разрушенная 
  придонная область удерживается силами адгезии и при задубливании оседает без искажения геометрии края.
\end{itemize}

\subsubsection{Химическая стойкость и адгезия}

\begin{itemize}
  \item \textbf{Химическая стойкость}: устойчивость защитного рельефа к кислотным и щелочным 
  травителям. Новолачная основа позитивных резистов устойчива к азотной кислоте ($48\%$-ной) 
  в течение $3 \dots 5$~мин при толщине $1$~мкм, но не обладает стойкостью к концентрированным щелочам.
  \item \textbf{Адгезия}: сцепление пленки фоторезиста с поверхностью подложки. Поскольку полимеры 
  фоторезиста гидрофобны, подложка также должна обладать гидрофобными свойствами. Свежеокисленная 
  поверхность кремния гидрофобна, однако на воздухе быстро сорбирует пары воды и становится гидрофильной, 
  что приводит к отслаиванию резиста и подтравливанию. Для восстановления адгезии пластины подвергают 
  высокотемпературной сушке ($700 \dots 800~^\circ\mathrm{C}$) либо химической гидрофобизации парами 
  гексаметилдисилазана (ГМДС).
\end{itemize}

\subsection{Технологический маршрут фотолитографии}

Процесс фотолитографии включает следующую строгую последовательность операций:
\begin{enumerate}
  \item \textbf{Контроль и подготовка подложки}: визуальный оптический контроль пластин, химическая 
  очистка от органических и механических загрязнений, гидрофобизация поверхности.
  \item \textbf{Нанесение фоторезиста}: дозированная подача раствора фоторезиста на пластину с 
  последующим центрифугированием со скоростью $2000 \dots 10000$~об/мин в течение $20 \dots 30$~с. 
  Толщина пленки определяется вязкостью раствора и скоростью вращения. Время между подачей резиста 
  и разгоном центрифуги минимизируют ($0{,}5 \dots 1$~с), чтобы исключить преждевременное испарение 
  растворителя. На периферии пластины формируется краевой валик утолщения.
  \item \textbf{Сушка фотослоя (первая термообработка)}: нагрев при температуре $90 \dots 100~^\circ\mathrm{C}$ 
  для удаления растворителя и релаксации внутренних механических напряжений в пленке. Наиболее 
  эффективен длинноволновый ИК-нагрев ($6 \dots 20$~мкм): излучение отражается от кремния и нагревает 
  резист снизу вверх, что исключает образование поверхностной корки и пузырей. Предел температуры 
  ограничен термическим распадом (термолизом) НХД.
  \item \textbf{Совмещение и экспонирование}:
  \begin{itemize}
    \item Шаблон ориентируют относительно базового среза пластины, затем проводят совмещение по 
    специальным реперным знакам (вписывание круга в круг, креста в крест) с точностью $\pm (0{,}2 \dots 0{,}5)$~мкм;
    \item Экспонирование актиничным излучением задает дозу $H = E \cdot t$, необходимую для фотолиза НХД.
  \end{itemize}
  \item \textbf{Постэкспозиционная термообработка (PEB)}: кратковременный отжиг пластины для устранения 
  неоднородности профиля резиста, вызванной стоячими световыми волнами при отражении от подложки.
  \item \textbf{Проявление}: выдержка пластины в щелочном проявителе ($\mathrm{KOH}$, тринатрийфосфат) 
  в течение $20 \dots 30$~с с последующей промывкой деионизованной водой и сушкой азотом. Качество 
  контролируется по величине клина проявления, не превышающего толщину резиста.
  \item \textbf{Задубливание (вторая термообработка)}: ступенчатый нагрев при $100 \pm 5~^\circ\mathrm{C}$ 
  (до $140~^\circ\mathrm{C}$). Удаляются остатки проявителя и воды, происходит термическая сшивка смолы, 
  слой уплотняется, существенно возрастают его адгезия и химическая стойкость.
  \item \textbf{Травление подслоя ($\mathrm{SiO}_2$)}: жидкостное химическое травление вскрытых окон 
  в буферном травителе (BHF):
  \begin{itemize}
    \item В чистом водном растворе плавиковой кислоты травление диоксида кремния сопровождается 
    выделением газообразного тетрафторида кремния:
    \begin{equation*}
      \mathrm{SiO}_2 + 4\mathrm{HF} \to \mathrm{SiF}_4\uparrow + 2\mathrm{H}_2\mathrm{O}.
    \end{equation*}
    Пузырьки газа отслаивают фоторезист от подложки, вызывая растравливание рисунка;
    \item Для подавления газовыделения используют буферный травитель — смесь $\mathrm{HF}$ и $\mathrm{NH}_4\mathrm{F}$. 
    Ионы фтора переводят кремний в растворимый жидкий гексафторосиликатный комплекс:
    \begin{equation*}
      \mathrm{SiF}_4 + 2\mathrm{F}^- \rightleftarrows [\mathrm{SiF}_6]^{2-}, \quad 
      2\mathrm{NH}_4^+ + [\mathrm{SiF}_6]^{2-} \rightleftarrows (\mathrm{NH}_4)_2\mathrm{SiF}_6.
    \end{equation*}
    Скорость травления $\mathrm{SiO}_2$ в BHF составляет около $0{,}08$~мкм/мин (при $20 \dots 22~^\circ\mathrm{C}$). 
    Вскрытые окна контролируют под микроскопом по серебристо-серому цвету кремния.
  \end{itemize}
  \item \textbf{Снятие фоторезиста}: деструкция полимерной маски в горячей серной кислоте 
  $\mathrm{H}_2\mathrm{SO}_4$, органических растворителях (диметилформамид, ацетон) либо 
  плазмохимическим озолением в кислородной ВЧ-плазме.
\end{enumerate}

\subsection{Влияние времени экспонирования на качество рельефа}

Точность передачи топологических размеров элементов критически зависит от времени экспонирования $t$:
\begin{enumerate}
  \item \textbf{Недоэкспонирование}: доза излучения недостаточна для полного разрушения молекул НХД по 
  всей глубине слоя. В проявителе экспонированные окна очищаются не полностью, остается тонкая пленка 
  резиста (пленка недопроявления), препятствующая травлению подслоя.
  \item \textbf{Переэкспонирование}: при избыточной дозе свет за счет дифракции на краях элементов шаблона 
  и рассеяния в толще пленки проникает в область геометрической тени под непрозрачные участки маски. 
  Это приводит к уширению вскрываемых окон в позитивном резисте и уменьшению ширины защитных линий.
\end{enumerate}

При контактной фотолитографии наличие микрозазора между шаблоном и пластиной вызывает дифракционное 
уширение пучка под углом $\theta_\text{диф} \propto \sqrt{\lambda / d_\text{зазор}}$, что приводит к нерезкости 
краев и возникновению клина растравливания. Поэтому для качественного воспроизведения микроструктур 
требуется прецизионный выбор оптимального времени экспозиции и минимизация рабочего зазора.

\end{document}

\end{document}